\clearpage
\fancyhead[L]{\small GAUSS AN VON ZACH.}
\hypertarget{native-286-lower}{}
\begin{center}
\textbf{GAUSS AN VON ZACH.}

\hrulefill

Monatliche Correspondenz. Band XV. S. 377. 378. 1807 April.

\hrulefill
\end{center}

\begin{flushright}
Braunschweig, den 10. März 1807.
\end{flushright}

\par\indent Hiebei habe ich die Ehre, Ihnen meine ersten Beobachtungen der \emph{Pallas} in diesem Jahre zu
übersenden. Das schlechte Wetter hat bisher die Fortsetzung gehindert.

\begin{center}\begingroup
\fontsize{9}{11}\selectfont
\setlength{\tabcolsep}{3pt}\renewcommand{\arraystretch}{1.15}
\begin{tabular}{l|l|l|l}
1807 & Mittl. Zeit & \shortstack{Scheinb. gerade\\Aufsteigung $\PallasSymbol$} & \shortstack{Scheinb. nördliche\\Abw. $\PallasSymbol$} \\
\hline
Februar 28 & $11^{\mathrm u}\,55^{\mathrm m}\,42^{\mathrm s}$ & $236^\circ\;48'\;6''$ & $8^\circ\;10'\;15''$ \\
März \phantom{0}3 & $12\;\;55\;\;25$ & $237\;\;15\;\;5$ & $8\;\;54\;\;19$ \\
--- \phantom{0}5 & $12\;\;22\;\;2$ & $237\;\;30\;\;41$ & $9\;\;24\;\;52$ \\
\end{tabular}\endgroup\end{center}

\par\indent Die beiden ersten Beobachtungen sind weniger zuverlässig, besonders in Ansehung der Declinationen; hingegen ist die dritte unter günstigern Umständen gemachte, so genau, als es meine Hülfs%
\SourceLeaf{287}%
mittel zulassen. Zur Vergleichung dienten am fünften März zwei Sterne neunter Grösse aus der
\emph{Hist. cél.}, deren scheinbare Positionen angenommen wurden wie folgt

\begin{center}\begingroup
\fontsize{9}{11}\selectfont
\setlength{\tabcolsep}{3pt}\renewcommand{\arraystretch}{1.15}
\begin{tabular}{l|l}
Gerade Aufsteig. & Nördl. Abweich. \\
\hline
$237^\circ\;19'\;12''$ & $9^\circ\;27'\;50''$ \\
$237\;\;31\;\;13$ & $9\;\;28\;\;8$ \\
\end{tabular}\endgroup\end{center}

\par\indent Meine Ephemeride gibt hienach die Declinationen genau, die geraden Aufsteigungen 12 bis 13
Min. zu gross; letztere Differenz scheint noch im Zunehmen zu sein. Für die Mailänder Beobachtungen
der \emph{Ceres} von 1806 danke ich verbindlichst; ich habe dadurch nun auch diese Opposition in gewünschter Schärfe. Nur finde ich, dass die Declination am 4. Januar um 1 Min. zu klein ist und $29^\circ\;21'\;29''3$
sein sollte (Februar-Heft 1807 S. 159). Dies ist kein Schreibfehler, denn auch die Breite ist zu klein,
und Herr \textsc{Carlini} muss sich zufälligerweise auch bei der aus den Elementen berechneten Breite um
1 Min. geirrt haben. Machen Sie diesen geschickten Astronomen doch gelegentlich aufmerksam darauf.
Mein letztes Resultat stimmt fast genau mit dem von Herrn \textsc{Carlini} berechneten überein. Ich finde die
Opposition den 3. Jan. 1806, $12^{\mathrm u}\,20^{\mathrm m}\,12^{\mathrm s}$ M.\,Z. in Mailand in $103^\circ\;1'\;9''4$ wahrer Länge und
$6^\circ\;28'\;32''4$ geoc. nördl. Breite. Mit Benutzung dieses Resultats werde ich nun warten, bis auch die
heurige Opposition beobachtet sein wird.

\clearpage\fancyhead[L]{\small VESTA.}
\begin{center}
\textbf{VESTA.}

\hrulefill

Monatliche Correspondenz. Band XV. S. 502–508. 1807 Mai.

\hrulefill
\end{center}

--- --- ---

\par\indent „Sobald ich mit einiger Gewissheit glauben konnte“, (schreibt uns Dr. \textsc{Olbers} unterm 22. April)
„wirklich der erste Entdecker dieses neuen Planeten zu sein, bat ich unsern unvergleichlichen \textsc{Gauss},
der sich so ausnehmende Verdienste um alle diese kleinen Planeten erworben hat, ihm Namen und
Zeichen zu bestimmen. Dr. \textsc{Gauss} hat meine Bitte erfüllt; `Sie legen mir, antwortete er, die Ehre, bei
Ihrem Planeten Pathenstelle zu vertreten, so dringend an's Herz, dass ich mich derselben nicht entziehen kann, so wenig ich auch Anspruch darauf habe. Es sei also darum; ich weiss dem Planeten keinen schönern Namen zu geben, als den der Göttin, die die Völker der alten Zeit zur Schutzgöttin der
reinen Sitten, der makellosen Tugend und des häuslichen Glückes machten. Finden Sie also meine
Wahl nicht unschicklich: so heisse Ihr Töchterchen: \emph{Vesta!}' Ich finde diesen Namen sehr glücklich
gewählt. Als Zeichen hat Herr Dr. \textsc{Gauss} die symbolische Vorstellung des auf dem Altare der Göttin
brennenden heiligen Feuers \VestaAltarBody{} bestimmt, und auch dies scheint mir in aller Absicht seinem Endzweck
zu entsprechen.“

\begin{flushright}
[\textsc{von Zach.}]
\end{flushright}
\clearpage
\fancyhead[L]{\small GAUSS AN BODE.}
\SourceLeaf{288}%
\begin{center}
\textbf{GAUSS AN BODE.}

\hrulefill

Astronomisches Jahrbuch für 1810. S. 210–214. Berlin 1807.

\hrulefill
\end{center}

\begin{flushright}
Braunschweig, den 8. Mai 1807.
\end{flushright}

\par\indent Ew. --- sage ich vielen Dank für die mir gütigst mitgetheilten Beobachtungen der Vesta. Ich
würde Ihnen die Resultate, die ich schon vor 18 Tagen gefunden hatte, geschickt haben, wenn ich
nicht anfangs Mistrauen in dieselben setzen, und hernach die Aussicht hätte haben müssen, sie bald
verbessern zu können. Das letztre ist auch jetzt noch nicht thunlich, wie Sie aus dem folgenden
sehen werden.

\par\indent Beobachtet habe ich die Vesta nur zwei Abende; die Beobachtungen wurden mir bald durch den
hohen Stand des Planeten erschwert, und ich konnte sie nachher um so eher aufgeben, da die Herren
\textsc{Olbers} und \textsc{Harding} die Güte gehabt haben, mir immer posttäglich ihre Beobachtungen zuzuschicken.
Hier sind meine beiden:

\begin{center}\begingroup
\fontsize{9}{11}\selectfont
\setlength{\tabcolsep}{3pt}\renewcommand{\arraystretch}{1.15}
\begin{tabular}{l|l|l}
& \multicolumn{2}{c}{Scheinbare} \\
& Ger. Aufst. \VestaAltarRA{} & Abweich. \VestaAltarDecl{} \\
\hline
1807. April 6. $8^{\mathrm u}\,22^{\mathrm m}\,45^{\mathrm s}$ M.\,Z. & $182^\circ\;21'\;54''6$ & $12^\circ\;28'\;0''9$ \\
\phantom{1807. April} 8. $\;\;8\;\;\;\;28\;\;\;\;49$ & $181\;\;\;56\;\;\;51.0$ & $12\;\;\;35\;\;\;45.8$ \\
\end{tabular}\endgroup\end{center}

\par\indent Am 20. April, als ich Dr. \textsc{Olbers} Beobachtung vom 17., und somit ein Zeitintervall von 19 Tagen
erhalten hatte, unternahm ich die erste Bestimmung der Bahn dieses neuen Planeten. Ich theilte
auch bereits am 21. die gefundenen Elemente Hrn. Dr. \textsc{Olbers} mit, konnte es aber bei der so kurzen
Zwischenzeit, der beschränkten Genauigkeit der Beobachtungen, die grösstentheils nur am Kreismikrometer angestellt waren, und dem hier eintretenden der genauen Bestimmung nachtheiligen
Umstande, dass die Neigung der Bahn nicht gross ist, noch nicht wagen, diese Resultate auch nur
für Annäherungen mit Zuverlässigkeit auszugeben. Ich behielt mir daher vor, sie zu verbessern,
sobald ich durch spätere Beobachtungen dazu in den Stand gesetzt würde. Allein die nachher
erhaltenen Beobachtungen stimmten fortwährend mit den Elementen so genau überein, dass sich der
Unterschied noch gar nicht ausmitteln liess, und da es also eigentlich eben so gut ist, als wenn die
Elemente unmittelbar mit auf die späteren Beobachtungen gegründet waren, so glaube ich jetzt, dass
man jene nun doch mit Zuversicht als erste Annäherung gelten lassen darf. Erst mit dem Ende des
Aprils und dem Anfange dieses Monats hat die Differenz angefangen merklich zu werden, und man
sieht nun wenigstens so viel, dass jetzt die berechneten Rectascensionen etwas zu gross, die Declinationen etwas zu klein sind, aber die eigentliche Grösse dieser Differenzen ist noch zu klein und
wird durch die Beobachtungen noch zu schwankend angedeutet, als dass jetzt schon eine zweite
Berechnung der Bahn mit Nutzen vorgenommen werden könnte. Sobald der Fehler etwas beträchtlicher geworden sein wird, soll dies geschehen. Hier also meine \emph{Elemente} I. \emph{der Vesta.} (Sie stehen
schon oben S. 198 dieses Jahrbuches, \textsc{Bode}.) [Sie sind in diesem Bande auf S. 286 abgedruckt.]
\SourceLeaf{289}%
\par\medskip\noindent Vergleichung mit sämmtlichen bis jetzt mir bekannt gewordenen Beobachtungen.

\begingroup\fontsize{9}{11}\selectfont\setlength{\tabcolsep}{2.5pt}\renewcommand{\arraystretch}{1.10}
\begin{longtable}{l|l|l|l|l|l|l}
1807 & Mittl. Zeit & Berechn. AR. & Differ. & Berechn. Decl. & Differ. & Beobachter \\
\hline
März 29 & $8^{\mathrm u}\,21^{\mathrm m}\,20^{\mathrm s}$ & $184^\circ\;8'\;54''4$ & $-1''4$ & . . . . & . . . & Olbers \\
        & $10\;\;31\;\;16$ & $184\;\;\;7\;\;40.2$ & $-10.6$ & $11^\circ\;47'\;52''4$ & $+5''4$ & Olbers \\
30      & $8\;\;34\;\;53$ & $183\;\;55\;\;\;3.0$ & $+6.2$ & . . . . & . . . & Olbers \\
        & $8\;\;44\;\;\;8$ & $183\;\;54\;\;57.7$ & $+11.9$ & $11\;\;53\;\;18.6$ & $+7.6$ & Olbers \\
        & $12\;\;33\;\;17$ & $183\;\;52\;\;44.2$ & $+3.4$ & $11\;\;54\;\;14.7$ & $-12.3$ & Olbers \\
April 1 & $9\;\;50\;\;\;0$ & $183\;\;27\;\;\;3.7$ & $+0.9$ & $12\;\;\;4\;\;40.7$ & $-11.3$ & Olbers \\
        & $12\;\;\;4\;\;35$ & $183\;\;25\;\;47.8$ & $+23.6$ & $12\;\;\;5\;\;\;9.6$ & $-22.1$ & Bessel \\
2       & $8\;\;21\;\;\;1$ & $183\;\;14\;\;23.8$ & $+6.0$ & $12\;\;\;9\;\;34.7$ & $-12.3$ & Olbers \\
        & $11\;\;38\;\;13$ & $183\;\;12\;\;33.7$ & $+19.6$ & $12\;\;10\;\;16.5$ & $-17.7$ & Bessel \\
3       & $8\;\;16\;\;49$ & $183\;\;\;1\;\;\;3.6$ & $+2.8$ & $12\;\;14\;\;33.6$ & $-6.4$ & Olbers \\
4       & $9\;\;\;9\;\;43$ & $182\;\;47\;\;20.9$ & $+0.1$ & $12\;\;19\;\;29.1$ & $-15.9$ & Olbers \\
5       & $11\;\;17\;\;\;3$ & $182\;\;33\;\;10.7$ & $-0.2$ & $12\;\;24\;\;21.5$ & $+2.6$ & Harding \\
6       & $8\;\;22\;\;45$ & $182\;\;21\;\;48.5$ & $-6.1$ & $12\;\;28\;\;\;6.7$ & $+5.6$ & Gauss \\
        & $11\;\;12\;\;15$ & $182\;\;20\;\;21.6$ & $-12.2$ & $12\;\;28\;\;34.2$ & $+45$ & Harding \\
8       & $8\;\;21\;\;37$ & $181\;\;56\;\;42.9$ & $-0.1$ & $12\;\;35\;\;50.1$ & $-8.0$ & Olbers \\
        & $8\;\;28\;\;49$ & $181\;\;56\;\;42.2$ & $-8.8$ & $12\;\;35\;\;50.1$ & $+4.3$ & Gauss \\
        & $11\;\;\;2\;\;43$ & $181\;\;55\;\;22.5$ & $+3.3$ & $12\;\;36\;\;13.6$ & $+1.4$ & Harding \\
9       & $8\;\;30\;\;46$ & $181\;\;44\;\;25.6$ & $+8.6$ & $12\;\;39\;\;19.8$ & $+5.7$ & Olbers \\
        & $10\;\;57\;\;58$ & $181\;\;43\;\;13.8$ & $+11.5$ & $12\;\;39\;\;39.6$ & $+3.5$ & Harding \\
12      & $8\;\;27\;\;20$ & $181\;\;\;9\;\;25.6$ & $-8.4$ & $12\;\;48\;\;\;7.4$ & $-8.7$ & Olbers \\
13      & $10\;\;39\;\;14$ & $180\;\;57\;\;28.3$ & $+11.3$ & $12\;\;50\;\;41.5$ & $-1.5$ & Bode \\
14      & $8\;\;19\;\;51$ & $180\;\;47\;\;36.3$ & $-21.7$ & $12\;\;52\;\;37.2$ & $-10.9$ & Olbers \\
        & $10\;\;34\;\;33$ & $180\;\;46\;\;39.2$ & $-0.7$ & $12\;\;52\;\;44.4$ & $+4.7$ & Harding \\
17      & $8\;\;12\;\;35$ & $180\;\;17\;\;21.6$ & $-5.2$ & $12\;\;57\;\;18.3$ & $+2.3$ & Olbers \\
23      & $8\;\;30\;\;52$ & $179\;\;26\;\;39.7$ & $-1.3$ & $12\;\;59\;\;20.4$ & $+0.4$ & Olbers \\
        & $9\;\;53\;\;49$ & $179\;\;26\;\;16.1$ & $-1.2$ & $12\;\;59\;\;18.9$ & $+4.6$ & Harding \\
24      & $9\;\;49\;\;26$ & $179\;\;19\;\;20.7$ & $+4.7$ & $12\;\;58\;\;42.5$ & $-12.5$ & Bode \\
25      & $8\;\;44\;\;38$ & $179\;\;13\;\;\;0.6$ & $-18.4$ & $12\;\;57\;\;52.4$ & $-5.6$ & Olbers \\
        & & $179\;\;12\;\;57.1$ & $+8.1$ & $12\;\;57\;\;51.9$ & $-14.1$ & Olbers \\
        & $9\;\;45\;\;\;5$ & $179\;\;12\;\;45.9$ & $+9.9$ & $12\;\;57\;\;50.1$ & $-16.9$ & Bode \\
        & $10\;\;25\;\;19$ & $179\;\;12\;\;28.7$ & $+5.6$ & $12\;\;57\;\;47.4$ & $+17.0$ & Bessel \\
26      & $8\;\;46\;\;\;4$ & $179\;\;\;6\;\;45.0$ & $-3.0$ & $12\;\;56\;\;44.2$ & $-6.8$ & Olbers \\
        & $9\;\;40\;\;45$ & $179\;\;\;6\;\;35.9$ & $-0.1$ & $12\;\;56\;\;42.3$ & $-15.7$ & Bode \\
        & $9\;\;40\;\;42$ & $179\;\;\;6\;\;32.5$ & $+12.6$ & $12\;\;56\;\;41.6$ & $-7.1$ & Harding \\
        & $10\;\;44\;\;12$ & $179\;\;\;6\;\;15.7$ & $+5.0$ & $12\;\;56\;\;38.0$ & $+8.5$ & Bessel \\
27      & $9\;\;36\;\;26$ & $179\;\;\;0\;\;51.3$ & $+3.3$ & $12\;\;55\;\;19.0$ & $-3.0$ & Bode \\
        & $9\;\;36\;\;23$ & $179\;\;\;0\;\;48.2$ & $+16.1$ & $12\;\;55\;\;18.1$ & $-10.6$ & Harding \\
        & $10\;\;45\;\;27$ & $179\;\;\;0\;\;31.3$ & $+7.3$ & $12\;\;55\;\;13.5$ & $+3.5$ & Bessel \\
        & $11\;\;18\;\;35$ & $179\;\;\;0\;\;23.6$ & $+12.6$ & $12\;\;55\;\;11.3$ & $-9.7$ & Olbers \\
28      & $10\;\;57\;\;\;6$ & $178\;\;55\;\;10.4$ & $+8.9$ & $12\;\;53\;\;32.9$ & $-22.3$ & Bessel \\
        & $11\;\;34\;\;20$ & $178\;\;55\;\;\;2.4$ & $+3.4$ & $12\;\;53\;\;30.1$ & $-20.9$ & Olbers \\
29      & $9\;\;27\;\;52$ & $178\;\;50\;\;38.9$ & $+17.9$ & $12\;\;51\;\;47.0$ & $-15.0$ & Bode \\
        & $9\;\;28\;\;51$ & $178\;\;50\;\;36.0$ & $-0.2$ & $12\;\;51\;\;45.8$ & $-2.6$ & Harding \\
        & $11\;\;\;4\;\;54$ & $178\;\;50\;\;16.6$ & $+21.8$ & $12\;\;51\;\;37.3$ & $-37.8$ & Bessel \\
30      & $9\;\;23\;\;38$ & $178\;\;46\;\;11.6$ & $+23.6$ & $12\;\;49\;\;38.7$ & $-21.3$ & Bode \\
        & $9\;\;23\;\;39$ & $178\;\;46\;\;\;9.1$ & $-13.1$ & $12\;\;49\;\;37.4$ & $+1.1$ & Harding \\
Mai 1   & $9\;\;19\;\;26$ & $178\;\;42,\;\;10.7$ & $+14.7$ & $12\;\;47\;\;15.7$ & $-15.3$ & Bode \\
        & $11\;\;23\;\;35$ & $178\;\;41\;\;48.2$ & $+22.2$ & $12\;\;47\;\;\;0.8$ & $-51.2$ & Olbers \\
2       & $9\;\;15\;\;14$ & $178\;\;38\;\;34\;3$ & $+24.1$ & $12\;\;44\;\;36.7$ & $-27.4$ & Harding \\
4       & $9\;\;\;7\;\;\;1$ & $178\;\;32\;\;46.7$ & $+17.7$ & $12\;\;38\;\;40.5$ & $-23.5$ & Bode \\
\end{longtable}\endgroup

\par\indent Ich freue mich, dass die von mir gewählte Benennung und Bezeichnung der \emph{Vesta} Ihren Beifall
erhalten hat. [Die mittleren Zeiten sind die der Beobachtungsorte.]
